\section*{الفصل \ref{ch:b2:chemical-potential} --- الجهد الكيميائي والخلائط}

\begin{solution}{exo:b2:chemical-potential:1}
يقطع المماس عند $x_2 = 0.4$ الحافتين عند $\bar V_1 \approx 14$ و$\bar
V_2 \approx 49$ (بوحدات النموذج). ومنه $0.6 \times 14 + 0.4 \times 49 = 28.0$،
وهو الحجم المولي المقروء على المنحنى عند $x_2 = 0.4$، وهو أدنى من
$0.6 \times 18 + 0.4 \times 58 = 34.0$ للحجوم النقية.
\end{solution}

\begin{solution}{exo:b2:chemical-potential:2}
$n(\text{ماء}) = 1000/18.015 = \qty{55.5}{mol}$، $x_1 = 55.5/56.5 = 0.982$؛
$p = 0.982 \times 3.17 = \qty{3.11}{kPa}$: انخفاض قدره 1.8\,\%.
\end{solution}

\begin{solution}{exo:b2:chemical-potential:3}
$p(\ce{O2}) = 0.2095 \times \qty{1.013e5}{Pa} = \qty{2.12e4}{Pa}$؛
$c = 1.3 \times 10^{-5} \times 2.12 \times 10^4 = \qty{0.28}{mol/m^3} =
\qty{2.8e-4}{mol/L}$، أي $2.8 \times 10^{-4} \times 32.0 \times 10^3 =
\qty{8.8}{mg/L}$.
\end{solution}

\begin{solution}{exo:b2:chemical-potential:4}
$Q = (1.0)^2/[(1.0)(0.010)^3] = 1.0 \times 10^6$؛ $RT\ln Q = 2.479 \times 13.82 =
\qty{34.2}{kJ/mol}$؛ $\Delta_r G = -32.8 + 34.2 = \qty{+1.4}{kJ/mol} > 0$: ففي هذا
الخليط الفقير جدًا بالهيدروجين تتفكك الأمونيا.
\end{solution}

\begin{solution}{exo:b2:chemical-potential:5}
جيبس--دوهيم: $x_1\dd\bar V_1 + x_2\dd\bar V_2 = 0$، ومنه $\dd\bar V_2 =
-(0.70/0.30) \times (-0.30) = \qty{+0.70}{mL/mol}$.
\end{solution}

\begin{solution}{exo:b2:chemical-potential:6}
$c(\text{جسيمات}) = 2 \times 9.0/58.44 = \qty{0.308}{mol/L} =
\qty{308}{mol/m^3}$؛ $\Pi = 308 \times 8.314 \times 310 = \qty{7.9e5}{Pa}$،
أي نحو \qty{7.9}{bar}. ولداخل كرية الدم الحمراء \omterm{def:b2:chemical-potential:osmotic-pressure}{الضغط الأسموزي} نفسه:
فلا تدفق صافيًا للماء. أما في الماء النقي فيدخل الماء الخلية، فتنتفخ
وتنفجر.
\end{solution}

\begin{solution}{exo:b2:chemical-potential:7}
$c = \Pi/RT = 825/(8.314 \times 298.15) = \qty{0.333}{mol/m^3}$؛ وفي
\qty{1.000e-4}{m^3}، $n = \qty{3.33e-5}{mol}$، $M = 2.00/3.33 \times 10^{-5}
= \qty{6.0e4}{g/mol}$ (\qty{60}{kg/mol}). $\Delta T_f = 1.86 \times 3.33
\times 10^{-5}/0.100 = \qty{6.2e-4}{K}$: لا يمكن قياسه، في حين يعادل \omterm{def:b2:chemical-potential:osmotic-pressure}{الضغط
الأسموزي} عمودًا من الماء طوله \qty{8}{cm}.
\end{solution}

\begin{solution}{exo:b2:chemical-potential:8}
(أ) $I = \qty{1.0e-3}{mol/kg}$، $\log\gamma_\pm = -0.51 \times 0.0316$،
$\gamma_\pm = 0.964$. (ب) $I = \frac12(4 \times 1.0 + 1 \times 2.0) \times
10^{-3} = \qty{3.0e-3}{mol/kg}$، $\log\gamma_\pm = -0.51 \times 2 \times 0.0548$،
$\gamma_\pm = 0.879$. (ج) $I = \frac12(4 + 4) \times 10^{-3} = \qty{4.0e-3}{mol/kg}$،
$\log\gamma_\pm = -0.51 \times 4 \times 0.0632$، $\gamma_\pm = 0.74$.
\end{solution}

\begin{solution}{exo:b2:chemical-potential:9}
$b = 0.310/1.86 = \qty{0.167}{mol/kg}$؛ $n = 0.167 \times 0.0500 =
\qty{8.33e-3}{mol}$؛ $M = 1.50/8.33 \times 10^{-3} = \qty{180}{g/mol}$ (سكر
سداسي، مثلًا).
\end{solution}

\begin{solution}{exo:b2:chemical-potential:10}
$x_1\,\dd(Ax_2^2) + x_2\,\dd(Ax_1^2) = 2Ax_1x_2(\dd x_2 + \dd x_1) = 0$ لأن
$x_1 + x_2 = 1$. وحين $x_1 \to 0$، يؤول $\gamma_1 \to \mathrm e^{A}$ ويكون $p_1 =
\gamma_1x_1p_1^* \approx \mathrm e^{A}p_1^*\,x_1$: $k_{H,1} = p_1^*\mathrm e^A$.
ومن أجل $A = 1.1$، $\mathrm e^{1.1} = 3.0$: فميل هنري ثلاثة أضعاف ميل راوول.
\end{solution}

\begin{solution}{exo:b2:chemical-potential:11}
(أ) المذاب المتطاير يساهم بضغط بخاره الخاص؛ فلا يعود البخار مذيبًا نقيًا
ويسقط الاستنتاج (بخار المذيب النقي في توازن مع المحلول). (ب) إذا دخل
المذاب في الجسم الصلب انخفض \omterm{def:b2:chemical-potential:chemical-potential}{الجهد الكيميائي} للجسم الصلب أيضًا، بمقدار
$RT\ln x_1^{\text{s}}$. وإذا كان للمحلول الصلب تركيب السائل نفسه، انخفض
الخطان بالمقدار نفسه ولم تتغير نقطة التجمد: فالانخفاض يتطلب مذابًا
تلفظه البلورة.
\end{solution}

\begin{solution}{exo:b2:chemical-potential:12}
قياس انخفاض درجة التجمد: $b_{\min} = 0.01/1.86 = \qty{5.4e-3}{mol/kg}$؛ وقياس ارتفاع درجة الغليان:
$0.01/0.513 = \qty{1.9e-2}{mol/kg}$. وعند $\qty{5.4e-3}{mol/L} =
\qty{5.4}{mol/m^3}$، $\Pi = 5.4 \times 8.314 \times 298 = \qty{1.3e4}{Pa}$،
أي أكثر من متر من الماء؛ ومقياس \omterm{def:b2:chemical-potential:osmotic-pressure}{الضغط الأسموزي} يقرأ بضعة باسكالات. فقياس \omterm{def:b2:chemical-potential:osmotic-pressure}{الضغط الأسموزي}
أشد الطرائق حساسية بفارق كبير، ولهذا يقيس الكتل المولية الكبيرة
للبوليمرات، التي تكون مولالياتها ضئيلة.
\end{solution}

\begin{solution}{pb:b2:chemical-potential:1}
\textbf{1.} $\mu_1 = \mu_1^*(T, p) + RT\ln x_1$.
\textbf{2.} الديول $30/62.07 = \qty{0.483}{mol}$، والماء $70/18.015 =
\qty{3.886}{mol}$؛ $x_1 = 3.886/4.369 = 0.889$.
\textbf{3.} $p = 0.889 \times 3.17 = \qty{2.82}{kPa}$.
\textbf{4.} $b = 0.483/0.070 = \qty{6.90}{mol/kg}$.
\textbf{5.} يغلي الديول عند \qty{197}{\celsius}: فضغط بخاره في درجة حرارة
الغرفة أدنى بكثير من ضغط بخار الماء.
\textbf{6.} $\mu_1^{\text{s}}(T_f) = \mu_1^{*\text{l}}(T_f) + RT_f\ln x_1$.
\textbf{7.} $\ln x_1 = -\Delta_{\text{انصهار}}G(T)/(RT)$ مع
$\Delta_{\text{انصهار}}G = \mu_1^{*\text{l}} - \mu_1^{\text{s}}$، المنعدم عند $T_f^*$؛
وعلاقة جيبس--هلمهولتز، $\dd(\Delta_{\text{انصهار}}G/T)/\dd T = -\Delta_{\text{انصهار}}H/T^2$،
بعد مكاملتها من $T_f^*$ إلى $T_f$، تعطي $\Delta_{\text{انصهار}}G(T_f)/T_f =
\Delta_{\text{انصهار}}H(1/T_f - 1/T_f^*)$، ومنها العلاقة.
\textbf{8.} $\ln x_1 \approx -x_2 \approx -bM_1$ و$1/T_f - 1/T_f^* \approx
-\Delta T_f/T_f^{*2}$: $\Delta T_f = (RT_f^{*2}M_1/\Delta_{\text{انصهار}}H)\,b$.
\textbf{9.} $333.4 \times 18.015 = \qty{6007}{J/mol}$؛ $K_f = 8.314 \times
273.15^2 \times 0.018015/6007 = \qty{1.86}{K.kg/mol}$.
\textbf{10.} أن $\Delta_{\text{انصهار}}H$ لا يتغير بين $T_f$ 
و$T_f^*$ (وهو يتغير، حسب \omterm{thm:b2:reaction-enthalpy:kirchhoff}{قانون كيرشوف}، بنحو 2\,\% لكل \qty{3}{K} هنا).
\textbf{11.} $\Delta T_f = 1.86 \times 6.90 = \qty{12.8}{K}$: \qty{-12.8}{\celsius}.
\textbf{12.} $1/T_f = 1/273.15 - (8.314/6007)\ln 0.889$: $T_f = \qty{261.6}{K}$،
أي \qty{-11.6}{\celsius}.
\textbf{13.} العلاقة المثالية الدقيقة، التي لا تفترض محلولًا مخففًا (فهنا
11\,\% من الجزيئات ديول). وكلتاهما تفترض \omterm{def:b2:chemical-potential:ideal-mixture}{خليطًا مثاليًا}
وأنتالبي انصهار ثابتًا؛ والماء والديول، المرتبطان بروابط هيدروجينية،
ليسا مثاليين، فتختلف نقطة التجمد الحقيقية.
\textbf{14.} جليد نقي: فالديول يبقى في السائل.
\textbf{15.} نزع الماء على هيئة جليد يخفض $x_1$ في السائل، وحسب
العلاقة تنخفض درجة حرارة التوازن: فيتجمد سائل التبريد على مجال
من درجات الحرارة، لا عند نقطة واحدة.
\textbf{16.} $K_b = 8.314 \times 373.12^2 \times 0.018015/40\,650 =
\qty{0.513}{K.kg/mol}$.
\textbf{17.} $\Delta T_b = 0.513 \times 6.90 = \qty{3.5}{K}$.
\textbf{18.} $1/T = 1/373.12 - (8.314/40\,650)\ln(2.0/1.013)$: $T = \qty{393.5}{K}$،
أي نحو \qty{120}{\celsius}.
\textbf{19.} الغطاء المضغوط: \qty{20}{K} مقابل \qty{3.5}{K} للديول.
\textbf{20.} $x_1 = \exp[-(6007/8.314)(1/253.15 - 1/273.15)] = 0.811$.
\textbf{21.} $x_2 = 0.189$: $w = 0.189 \times 62.07/(0.189 \times 62.07 + 0.811
\times 18.015) = 0.44$.
\textbf{22.} القانون المخفف: $b = 20/1.86 = \qty{10.75}{mol/kg}$، $w = 10.75 \times
62.07/(1000 + 10.75 \times 62.07) = 0.40$. فالتقريبان الخطيان ($\ln x_1 \approx
-x_2$، $x_2 \approx bM_1$) يبالغان في الانخفاض لكل مول من الديول عند تركيز
بهذا الارتفاع، فيطلب القانون المخفف ديولًا أقل.
\textbf{23.} في نموذج \omterm{def:b2:chemical-potential:ideal-mixture}{الخليط المثالي}، يبقى سائل التبريد ذو الكسر الكتلي من
الديول $\boldsymbol{w \approx 0.44}$ سائلًا حتى \qty{-20}{\celsius}.
\end{solution}
